\section{Introduction}

\begin{figure*}[!t]
\centering
\includegraphics[width=0.82\textwidth]{figures/fig_teaser.png}
\caption{Real-frame comparison and TALH inference. (a) Content-aligned
Native-HR (estimated), Trilinear, CLL-only, and TALH-Q outputs for
prompt 05, seed 9705, frame 41, with identical \(1.5\times\) crops. The
reference follows a complete 368p trajectory; handoff outputs are 720p.
(b) Native-HR performs all 50 evaluations at high resolution, whereas
TALH-Q and TALH-E execute the first 40 and 45 evaluations at low
resolution and retain 10 and 5 high-resolution suffix evaluations.}
\label{fig:teaser}
\end{figure*}

Recent latent video diffusion models such as Wan generate semantically
coherent videos with realistic motion and fine details \citep{wan2025},
but their inference cost grows aggressively with spatial resolution,
video length, and sampling steps. In DiT-based models, every denoising
iteration processes the full spatiotemporal token sequence. Executing the
entire trajectory at the target resolution forces early steps---which mainly
establish coarse structure and macro-motion---to bear the full burden of
high-resolution computation, severely limiting real-time deployment.
Since diffusion sampling inherently distributes generation across stages
(global layout in early steps, high-frequency textures in later steps),
a promising paradigm is mixed-resolution inference: modeling global
content in a low-resolution latent space first, then switching to the
target resolution for a short refinement suffix. This spatial-scaling
approach \citep{lightx2v2026,zheng2026mrflow} is fundamentally orthogonal
and complementary to temporal step-distillation methods.

Yet, dynamic resolution handoff is far more complex than simple tensor resizing,
collectively dictated by three error sources: trajectory gap, scale discrepancy,
and refinement residual. First, a substantial residual denoising gap persists
between the single-step prediction and the completed low-resolution endpoint.
Second, video VAE latents lack strict scale equivariance, failing to capture
the intricate cross-resolution correspondences jointly determined by the VAE
architecture and scale transitions. Finally, the lifted latents must seamlessly
re-enter the high-resolution trajectory for texture refinement.
Directly mapping the handoff prediction to the high-resolution endpoint forces a
single Joint Trajectory-Scale Lifter (JTSL) to resolve these challenges jointly.
Because the trajectory gap fluctuates with timesteps, schedulers, and prompts,
and since high-resolution details are underdetermined by low-resolution inputs,
such joint regression heavily suffers from high-frequency attenuation. Our
handoff-step scans confirm this, motivating a decomposition into localized subproblems.

We therefore propose \textbf{TALH} (Trajectory-Aligned Latent Handoff),
comprising three modular components activated at the handoff step: a LoRA-based
\textbf{Trajectory Alignment Adapter (TAA)} to calibrate the residual trajectory gap;
a \textbf{Clean Latent Lifter (CLL)} to map cross-scale clean latents; and a
\textbf{High-Resolution Trajectory Re-entry (HTR)} protocol for generative refinement.
Both TAA and CLL are trained via a model-endogenous self-distillation framework
supervised directly by the frozen target generator.

Our primary contributions are summarized as follows:
\begin{itemize}
\item \textbf{We propose TALH}, a lightweight mixed-resolution framework that shifts the bulk of denoising computations to a low-resolution prefix, providing controllable quality-efficiency Pareto operating points (TALH-Q/E) for online inference.
\item \textbf{We decouple dynamic resolution handoff} into three explicit phases via TAA, CLL, and HTR. This mathematical decomposition significantly mitigates the epistemic uncertainty inherent in joint trajectory-scale mapping, enabling isolated diagnostics.
\item \textbf{We construct a model-endogenous training paradigm} requiring no external data or teachers, and establish the consistency condition for cached-prefix states. Experiments on standard and distilled Wan samplers validate its robust acceleration and schedule compatibility.
\end{itemize}

\section{Related Work}

\subsection{Latent Video Generation and Cascaded Upscaling}
Latent Diffusion Models compress the generative process into a compact space for scalable high-resolution video synthesis \citep{rombach2022ldm}. Traditional frameworks or recent systems like LTX-2 and LUVE typically employ cascaded pipelines, using independent spatial latent upsamplers or post-generation refinement modules to elevate resolution \citep{blattmann2023videoldm, hacohen2026ltx2, zhao2026luve}. Similarly, video super-resolution methods focus on restoring already decoded RGB outputs \citep{zhou2024upscale, xu2024videogigagan}.

Unlike these approaches that treat upscaling as an isolated post-processing stage, TALH natively embeds the Clean Latent Lifter (CLL) into the original Wan sampling trajectory. By linking the lifted latents to subsequent steps via High-Resolution Trajectory Re-entry (HTR), cross-resolution scaling and generative priors are unified within a single, continuous inference process.

\subsection{Mixed-Resolution Diffusion Inference}
The changing-resolution inference in LightX2V utilizes fixed interpolation to adjust latent sizes mid-sampling \citep{lightx2v2026}. Concurrently, RALU demonstrates that training-free latent upscaling suffers from high-frequency aliasing and a resolution-dependent noise--timestep mismatch, designing a mixed-resolution scheme for Diffusion Transformers \citep{jeong2025ralu}. For image flow matching, MrFlow combines low-resolution structure generation with low-noise high-resolution refinement \citep{zheng2026mrflow}. TALH extends this paradigm to video trajectories by learning cross-resolution lifting directly within the model's latent space, while explicitly rectifying low-resolution tail-trajectory errors prior to the handoff.

\subsection{Few-Step Distillation and On-Policy Trajectory Learning}
To accelerate inference, temporal distillation methods minimize sampling steps via trajectory regression or consistency mapping \citep{luo2023lcm, yin2024dmd2, zhai2024mcm}. Recent works like Self-Forcing and AnyFlow further emphasize enforcing strict consistency between student training states and actual inference trajectories to mitigate distribution shift \citep{huang2025selfforcing, gu2026anyflow}.

TALH focuses on spatial rather than temporal computation, making it orthogonal to step reduction. Since TAA remains inactive prior to the handoff step, its parameter updates do not alter the low-resolution prefix trajectory. Leveraging this local activation property, Section~\ref{sec:cached} establishes the consistency conditions that allow cached prefix states to be directly accessed during inference without training--inference distribution drift.

\section{Method}

\subsection{Problem Definition and Overall Framework}

Let \(v_\phi(x_s, s, c)\) denote a frozen Wan flow predictor, where
\(x_s\) represents the intermediate latent before the \(s\)-th denoising
calculation, \(c\) is the text condition, and \(T\) is the total number
of denoising steps. We use the superscripts \(L\) and \(H\) to signify
the low-resolution and high-resolution latent spaces, respectively. Let
\(z_T^L\) denote the endpoint of the complete low-resolution trajectory
generated by the frozen model under a given text prompt, random seed,
initial noise, and low-resolution scheduler. We define
\((z_0^L, z_0^H)\) as the cross-resolution clean latent pair obtained by
encoding the low- and high-resolution versions of the same video through
the Wan VAE.

TALH introduces two independently trained modules: the Trajectory
Alignment Adapter (TAA) parameterized by LoRA weights
\(\Delta\theta_s\), and the Clean Latent Lifter (CLL) denoted as
\(U_\psi^{\mathrm{CLL}}\) with parameters \(\psi\). High-Resolution
Trajectory Re-entry (HTR) is a parameter-free re-routing operation. As
illustrated in the upper panel of Figure~\ref{fig:overview}, the complete data flow is
formulated as:

\[
x_s^L \xrightarrow{\mathrm{TAA}:v_{\phi+\Delta\theta_s}}
\widetilde z_s^L \xrightarrow{U_\psi^{\mathrm{CLL}}}
\widehat z_s^H \xrightarrow{\mathrm{HTR}} x_{s+1}^H.
\]

Crucially, TAA is dedicated solely to aligning the current
low-resolution clean prediction with the complete trajectory's
low-resolution endpoint, while CLL focuses exclusively on the
cross-resolution mapping of clean latents. These two modules are trained
separately without joint backpropagation. This decoupled training
strategy prevents the massive memory overhead of hosting both the 14B
denoiser and the CLL module simultaneously in the compute graph, while
enabling an isolated experimental analysis of trajectory and scale
errors.

The handoff step simultaneously governs three distinct types of errors.
Let \(E_{\mathrm{traj}}(s)\) denote the residual trajectory gap before
the handoff, \(E_{\mathrm{lift}}(s)\) be the cross-resolution lifting
error, and \(E_{\mathrm{refine}}(s)\) represent the remaining
high-resolution error left uncorrected after HTR. The total handoff
error can be conceptualized as:

\[
E_{\mathrm{handoff}}(s) \approx E_{\mathrm{traj}}(s)
+ E_{\mathrm{lift}}(s) + E_{\mathrm{refine}}(s).
\]

This conceptual decomposition reveals the central trade-off across
choices of \(s\). A later handoff reduces the residual trajectory gap
\(E_{\mathrm{traj}}(s)\) and moves the CLL input closer to its
clean-latent training distribution, thereby reducing
\(E_{\mathrm{lift}}(s)\). However, a later handoff also leaves fewer
high-resolution suffix evaluations and can therefore increase the
uncorrected refinement error \(E_{\mathrm{refine}}(s)\). TAA, CLL, and
HTR are designed to address these three error sources separately.

\begin{figure*}[!t]
\centering
\includegraphics[width=0.84\textwidth]{figures/fig_talh_overview.png}
\caption{Model-endogenous supervision for TAA and CLL. TAA pairs the
cached handoff state with the complete low-resolution trajectory endpoint
under identical prompt, seed, scheduler, and CFG settings. CLL pairs
clean latents encoded from a frozen-Wan-generated high-resolution video
and its RGB-downsampled counterpart. Both training-pair streams are
generated by the same frozen Wan model.}
\label{fig:overview}
\end{figure*}

\subsection{Model-Endogenous Dual-Supervision Data}

TALH operates entirely without external paired videos, pre-trained
super-resolution weights, or additional teacher models. The supervision
for both CLL and TAA is generated natively by the frozen target
generator, tracking two distinct conditional distributions as shown in
the lower panel of Figure~\ref{fig:overview}.

\textbf{Cross-Resolution Lifting Pairs.} We run the frozen Wan model
using a set of predefined text prompts and fixed random seeds to
generate high-resolution reference videos. These videos are spatially
downsampled in the RGB pixel space. Both the high- and low-resolution
variants are then encoded using the same Wan VAE to produce perfectly
aligned cross-resolution clean latent pairs \((z_0^L, z_0^H)\). This
construction preserves the frame count, semantic content, and motion
trajectories, forcing CLL to learn cross-scale correspondences tightly
bound to the target model's specific VAE and generative distribution.

\textbf{Trajectory Alignment Pairs.} Under identical text prompts,
random seeds, initial noise, and schedulers, we roll out the complete
low-resolution trajectory of the frozen model. We cache the intermediate
low-resolution state \(x_s^L\) immediately preceding the handoff step
and utilize the final low-resolution trajectory endpoint \(z_T^L\) as
the ground-truth target. Consequently, CLL receives cross-resolution
structural supervision, while TAA receives cross-timestep trajectory
supervision, forming a unified model-endogenous resolution-trajectory
self-distillation framework.

Although both types of supervision share the same base generative model,
they map to different conditional distributions: the low-resolution
input for CLL is derived from downsampled high-resolution video
encodings, whereas TAA directly targets the endpoints of native
low-resolution trajectories. Section~\ref{sec:factorialsetup} evaluates their end-to-end
synergy via a \(2\times2\) factorial design crossing trajectory states
and resolution operators.

\subsection{Cross-Resolution Clean Latent Lifting}

The Clean Latent Lifter (CLL) takes a 16-channel Wan VAE latent as input
and produces an upscaled 16-channel latent, preserving the temporal
dimension:

\[
U_\psi^{\mathrm{CLL}}:\mathbb{R}^{B\times16\times F\times h\times w}
\rightarrow \mathbb{R}^{B\times16\times F\times H\times W}.
\]

The architecture adopts an encoder-resampler-decoder structure inspired
by LTX-2. A 3D convolutional input layer projects the latents into a
256-dimensional hidden space, followed by 4 spatiotemporal residual
blocks to extract rich spatiotemporal features. A spatial resampler then
scales the resolution, after which 4 additional 3D residual blocks and a
final 3D convolutional layer reconstruct the 16-channel output. This
spatiotemporal processing enables the model to exploit information from
adjacent latent frames while keeping the total frame count invariant.

We implement two distinct spatial resampling configurations. For the
\(480\times832 \rightarrow 720\times1248\) transition, the latent grid
expands from \(60\times104\) to \(90\times156\). The network first
inflates the channel dimension by a factor of 9, applies a \(3\times\)
spatial PixelShuffle, and finally performs a stride-2 downsampling using
a fixed binomial blur kernel to achieve a rational \(3/2\times\)
upscaling:

\[
60\times104 \xrightarrow{\times3} 180\times312
\xrightarrow{/2} 90\times156.
\]

For the \(368\times640 \rightarrow 720\times1248\) transition, the
latent grid expands from \(46\times80\) to \(90\times156\). Since the
target size is not an exact integer multiple, the network executes a
\(2\times\) PixelShuffle to reach \(92\times160\), followed by a center
crop to isolate the exact \(90\times156\) target grid.

The training objective for CLL is defined as:

\[
\mathcal L_{\mathrm{CLL}}=\mathcal L_{\mathrm{rec}}
+\lambda_{\mathrm{content}}\mathcal L_{\mathrm{content}}
+\lambda_{\mathrm{temp}}\mathcal L_{\mathrm{temp}}.
\]

where

\[
\mathcal L_{\mathrm{rec}}
=\mathbb E\left[\sqrt{(\widehat z_0^H-z_0^H)^2+\epsilon^2}\right].
\]

\begin{align}
\mathcal L_{\mathrm{content}}
&=\lVert D(\widehat z_0^H)-z_0^L\rVert_1,\\
\mathcal L_{\mathrm{temp}}
&=\lVert\Delta_F\widehat z_0^H-\Delta_Fz_0^H\rVert_1.
\end{align}

Here, \(D\) denotes a downscaling operator that maps the predicted
high-resolution latent back to the low-resolution grid, and \(\Delta_F\)
represents the first-order temporal difference between adjacent latent
frames. We set \(\lambda_{\mathrm{content}}=0.2\),
\(\lambda_{\mathrm{temp}}=0.1\), and \(\epsilon=10^{-3}\) in our
experiments. These three loss terms jointly constrain point-wise
reconstruction, low-frequency content preservation, and temporal
consistency. Because a low-resolution latent cannot uniquely determine
high-resolution details, CLL focuses on recovering a stable cross-scale
representation, leaving the remaining high-frequency details to be
synthesized and corrected by the subsequent HTR and high-resolution
diffusion priors.

\subsection{Handoff-Local Trajectory Alignment}

TAA is implemented via LoRA \citep{hu2022lora} and is activated
exclusively at the designated handoff step. The base Wan model remains
completely frozen; low-rank updates of rank \(r\) are injected solely
into the \(q, k, v, o\) projections of the attention modules and the
\(ffn.0, ffn.2\) linear layers of the Feed-Forward Networks (FFNs):

\[
W'=W+\frac{\alpha}{r}BA.
\]

Under the flow matching parameterization, the TAA-aligned clean latent
prediction at step \(s\) is expressed as:

\[
\widetilde z_s^L
=x_s^L-\sigma_s v_{\phi+\Delta\theta_s}(x_s^L,s,c).
\]

The alignment target is the low-resolution trajectory endpoint \(z_T^L\)
obtained by fully running the frozen model up to step \(T\) under
identical conditions (text prompt, random seed, initial noise, and
scheduler). The trajectory alignment loss is formulated as:

\[
\mathcal L_{\mathrm{align}}
=\lVert\widetilde z_s^L-z_T^L\rVert_1
+0.1\lVert\widetilde z_s^L-z_T^L\rVert_2^2.
\]

For earlier handoffs (such as step 40), we introduce an additional
temporal residual alignment loss with a weight of 0.05 to mitigate the
propagation of frame-to-frame errors into the CLL module. While our
framework also supports direct supervision via the velocity target
\(v^*=\frac{x_s^L-z_T^L}{\sigma_s}\), all experiments in this paper
report results optimized using the clean latent alignment objective
above.

By compressing the residual error of the low-resolution tail trajectory
at the handoff step, TAA effectively projects the input of CLL back onto
the clean latent training domain. While the system-level acceleration of
TALH stems from the low-resolution prefix computations, TAA ensures
trajectory consistency at fixed handoff points. Earlier handoffs, which
inherently suffer from larger residual trajectory gaps, consequently
yield a broader optimization space for TAA calibration.

\subsection{Cached-Prefix State Consistency}
\label{sec:cached}

Assuming TAA is activated strictly at the handoff step \(s\), its LoRA
parameters satisfy \(\Delta\theta_k=0\) for all steps \(k < s\). Given a
deterministic sampler alongside fixed text embeddings, random seeds,
schedulers, and Classifier-Free Guidance (CFG) scales, the adapted
generative process shares the identical transition functions as the base
process prior to step \(s\). Thus, by induction over the sampling steps,
we establish:

\[
x_k^{\mathrm{adapted}}=x_k^{\mathrm{base}},
\qquad \forall k\leq s.
\]

This property guarantees that the pre-cached base model prefix state
\(x_s^L\) is mathematically identical to the state actually accessed by
TAA during inference. For localized adaptations restricted to a single
handoff step, caching the prefix state allows us to bypass redundant
prefix computations entirely during inference without introducing any
training-inference distribution drift, maintaining an \(O(1)\) state
retrieval overhead.

This consistency property relies on four strict conditions: TAA must be
engaged exclusively at the handoff step; training and inference must
share identical schedulers, CFG scales, and conditional inputs; the
prefix sampling path must be entirely deterministic; and no unaligned
stochastic operations can occur within the prefix trajectory. If TAA
were extended across multiple steps, its parameters would alter
subsequent states, requiring advanced training mechanisms such as
student model rollouts, synchronized teacher predictions, or Exponential
Moving Average (EMA) teacher models.

\subsection{High-Resolution Trajectory Re-entry and Two Instances}

Once TAA produces the aligned low-resolution clean latent, CLL outputs
the upscaled high-resolution clean latent:

\[
\widehat z_s^H=U_\psi^{\mathrm{CLL}}(\widetilde z_s^L).
\]

The High-Resolution Trajectory Re-entry (HTR) phase then introduces
target-resolution noise according to the schedule coefficient
\(\sigma_{s+1}\) of the next timestep:

\[
x_{s+1}^H
=(1-\sigma_{s+1})\widehat z_s^H
+\sigma_{s+1}\epsilon^H.
\]

The remaining denoising calculations from step \(s+1\) to \(T\) are
subsequently completed by the frozen Wan denoiser operating on the
high-resolution latent space. Unlike an \emph{Interpolated-Flow
Re-entry} strategy that directly interpolates low-resolution velocity
flows into the high-resolution space, TALH utilizes
\emph{Target-Resolution Noise Re-entry}. This design ensures that newly
introduced spatial frequencies are synthesized natively by the
high-resolution generative priors of the base model.

Our \textbf{50-step pipeline} uses Wan2.1 T2V-1.3B, configured with 50
inference steps, a sample shift of 8, and a CFG scale of 6. Extensive
handoff-step scanning identified two representative operational points.
The quality-centric \textbf{TALH-Q} (TALH@40) reserves 10
high-resolution steps, granting a larger budget for texture generation
while managing a wider residual trajectory gap that requires substantial
TAA alignment. Conversely, the efficiency-centric \textbf{TALH-E}
(TALH@45) reserves only 5 high-resolution steps to minimize latency; its
handoff state sits closer to a fully converged clean latent, resulting
in more stable CLL inputs and narrower TAA corrections. Together, these
two instances define the quality-efficiency Pareto frontier of TALH.

Our \textbf{4-step pipeline} leverages Wan2.1 T2V-14B
StepDistill-CfgDistill, operating on four nominal timesteps
\([1000, 750, 500, 250]\) with a sample shift of 5. The distilled
instance \textbf{TALH-D4} (TALH@3-of-4) runs low-resolution sampling for
the first two steps. At step 3, TAA predicts the aligned low-resolution
clean latent, which is scaled by CLL and injected back into the
high-resolution trajectory via HTR at step 4. A single final
high-resolution denoising step completes the generation.

\section{Experiments}

\subsection{Experimental Setup}

We evaluate TALH at both system and module levels, covering end-to-end
quality--latency trade-offs, the individual and interactive effects of
CLL and TAA, and compatibility with a 4-step distilled sampler.

For the 50-step pipeline, we generate approximately 1,000 reference
videos using Wan2.1 T2V-1.3B. Each video contains 81 frames at
\(720\times1248\) resolution and 16 fps. Bicubic RGB downsampling
produces corresponding \(480\times832\) and \(368\times640\) videos,
which are re-encoded by the frozen Wan VAE to construct aligned
cross-resolution latent pairs. TAA training pairs consist of the
low-resolution state immediately preceding a designated handoff and the
endpoint of the corresponding complete low-resolution trajectory.

The 4-step experiments use approximately 5,000 videos generated by
Wan2.1 T2V-14B StepDistill-CfgDistill. CLL is trained on
\(368\times640\rightarrow720\times1248\) clean latent pairs, while TAA
pairs the state preceding step 3 with the complete step-4
low-resolution endpoint. Training, validation, and test partitions are
separated by both text prompts and random seeds.

CLL is optimized with AdamW for up to 50k steps using a learning rate of
\(1\times10^{-4}\), \(bf16\) precision, and an EMA decay of 0.9999.
TAA is trained for up to 10k steps with AdamW, a learning rate of
\(5\times10^{-5}\), and LoRA rank/scale \(32/32\). Training uses four
NVIDIA H100 GPUs and requires approximately 8 hours for CLL and 33 hours
for TAA. All efficiency measurements are averaged over five independent
cold-start runs. Since all supervision is generated by the frozen target
model, TALH requires neither external paired videos nor an additional
upscaling teacher.

\subsection{Baselines and Evaluation Protocol}
\label{sec:factorialsetup}

System-level comparisons include \textbf{Native-HR Sampling}, which
executes all denoising steps at \(720\text{p}\);
\textbf{Trilinear Handoff}, which performs mixed-resolution inference
using fixed trilinear latent interpolation; and
\textbf{Endpoint Re-entry}, which completes the low-resolution
trajectory, applies CLL, and performs one low-intensity
high-resolution refinement step. We additionally report
\textbf{Native-LR Sampling} as a low-resolution reference.

Module interactions are evaluated through a \(2\times2\) factorial
design that crosses trajectory states
\(\{\text{Unaligned},\text{TAA-aligned}\}\) with resolution operators
\(\{\text{Trilinear},\text{CLL}\}\). The four resulting configurations
are Trilinear Handoff, TAA+Trilinear, CLL-only, and the complete TALH
framework. Oracle-Aligned CLL and Joint Trajectory-Scale Lifting (JTSL)
are used as diagnostic references; extended results are provided in the
Supplementary Material.

For trajectory alignment, we report endpoint latent \(L_1\), decoded
PSNR and LPIPS, temporal latent differences, and sample-wise win rates.
CLL is evaluated using latent \(L_1\), PSNR, SSIM, LPIPS, and temporal
\(L_1\). End-to-end quality is measured with VBench
\citep{huang2024vbench} and human evaluation. We define
\textbf{VBench-5} as the unweighted mean of subject consistency,
background consistency, motion smoothness, aesthetic quality, and
imaging quality.

All paired configurations share identical prompts, random seeds,
initial noise, schedulers, and guidance parameters. Prompt-level score
differences are analyzed using 10,000 paired bootstrap iterations for
\(95\%\) confidence intervals and two-sided exact sign tests. Human
evaluation uses randomized, double-blind comparisons over 10 matched
prompts, with three independent judges per prompt and prompt-level
majority voting. Inter-rater agreement is reported in the Supplementary
Material.

Efficiency measurements include end-to-end cold-start latency, peak GPU
memory, low- and high-resolution denoising evaluations, and speedup over
Native-HR Sampling. All TAA, CLL, HTR, VAE, and dynamic model-loading
overheads are included.

\begin{figure*}[!t]
\centering
\includegraphics[width=0.76\textwidth]{figures/fig_quality_efficiency.png}
\caption{Quality--efficiency trade-offs relative to Native-HR Sampling.
(a) TALH-Q and TALH-E provide quality- and efficiency-oriented operating
points, while Endpoint Re-entry achieves higher speed at a larger
quality cost. (b) Absolute per-dimension VBench changes relative to
Native-HR.}
\label{fig:quality}
\end{figure*}

\begin{table*}[!t]
\centering
\footnotesize
\setlength{\tabcolsep}{5pt}
\begin{tabular}{lrrrrrr}
\toprule
Method & LR/HR evals. & Time (s) $\downarrow$ & Latency red. &
Speedup $\uparrow$ & VBench-5 $\uparrow$ & Retention $\uparrow$ \\
\midrule
Native-HR Sampling & 0/50 & $253.10\pm1.53$ & --- &
$1.00\times$ & \textbf{0.82836} & 100.00\% \\
TALH-Q (Step 40) & 40/10 & $138.36\pm0.50$ & 45.33\% &
$1.83\times$ & 0.80983 & 97.76\% \\
TALH-E (Step 45) & 45/5 & $\mathbf{114.26\pm1.52}$ &
\textbf{54.85\%} & $\mathbf{2.22\times}$ & 0.80792 & 97.53\% \\
\bottomrule
\end{tabular}
\caption{System-level comparison against Native-HR Sampling. Timings
average five cold-start runs, and quality is evaluated on ten matched
prompts with all processing overheads included.}
\label{tab:system}
\end{table*}

\section{Results and Analysis}

\subsection{System-Level Quality--Efficiency Trade-off}

Table~\ref{tab:system} compares the standard 50-step Native-HR pipeline
with the two principal TALH operating points. TALH-Q moves the first 40
evaluations to the low-resolution grid and retains 10 high-resolution
suffix evaluations, reducing latency from 253.10 to 138.36 seconds
(\(1.83\times\)). TALH-E retains only five high-resolution evaluations
and further reduces latency to 114.26 seconds (\(2.22\times\)).
Their VBench-5 scores are 0.80983 and 0.80792, compared with 0.82836 for
Native-HR Sampling.

The Endpoint Re-entry baseline reaches 86.45 seconds
(\(2.93\times\)) but decreases VBench-5 to 0.80093, with inferior
scores on \(9/10\) prompts relative to Native-HR
(\(p=0.02148\)). TALH therefore retains a longer high-resolution suffix
when additional structural and texture refinement is beneficial, while
later handoffs prioritize inference speed. Peak memory remains
approximately 26 GiB across Native-HR, TALH-Q, TALH-E, and Endpoint
Re-entry, indicating that the primary system benefit of TALH is reduced
runtime rather than memory consumption.

\subsection{Component Analysis}

\paragraph{Clean latent lifting.}
Table~\ref{tab:operator} evaluates CLL against trilinear interpolation
over 50 paired clean latent samples. For
\(480\text{p}\rightarrow720\text{p}\), CLL reduces latent \(L_1\) from
0.1929 to 0.0991, corresponding to a \(48.6\%\) reduction. It also
improves decoded PSNR by 5.71 dB, reduces LPIPS from 0.2358 to 0.0629,
and lowers temporal \(L_1\) by \(30.8\%\). CLL wins on all 50 samples
for PSNR, LPIPS, and temporal consistency, and on 49/50 samples for
SSIM.

The more aggressive \(368\text{p}\rightarrow720\text{p}\) transition
shows the same trend: latent \(L_1\), LPIPS, and temporal \(L_1\) are
reduced by \(33.4\%\), \(51.1\%\), and \(15.5\%\), respectively.
These consistent latent, perceptual, and temporal improvements indicate
that CLL learns a substantially more faithful cross-resolution mapping
than fixed interpolation.

\begin{table*}[!t]
\centering
\footnotesize
\setlength{\tabcolsep}{5pt}
\begin{tabular}{llrrrrr}
\toprule
Input $\rightarrow$ Output & Operator & Latent $L_1\downarrow$ &
PSNR $\uparrow$ & SSIM $\uparrow$ & LPIPS $\downarrow$ &
Temp. $L_1\downarrow$ \\
\midrule
480p $\rightarrow$ 720p & Trilinear & 0.1929 & 24.29 & 0.7159 &
0.2358 & 0.02121 \\
480p $\rightarrow$ 720p & CLL & \textbf{0.0991} & \textbf{30.00} &
\textbf{0.8696} & \textbf{0.0629} & \textbf{0.01468} \\
368p $\rightarrow$ 720p & Trilinear & 0.2444 & 23.42 & 0.6740 &
0.3352 & 0.02388 \\
368p $\rightarrow$ 720p & CLL & \textbf{0.1628} & \textbf{24.10} &
\textbf{0.7592} & \textbf{0.1640} & \textbf{0.02017} \\
\bottomrule
\end{tabular}
\caption{CLL versus trilinear interpolation on paired clean latents
(\(n=50\)).}
\label{tab:operator}
\end{table*}

\paragraph{Trajectory alignment.}
At an alignment intensity of 0.75, TAA reduces the residual
low-resolution endpoint \(L_1\) error by \(25.82\%\) at step 40 and
\(21.03\%\) at step 45. Both settings improve all 10 evaluation
samples, with paired confidence intervals excluding zero. TAA also
increases decoded PSNR by 2.55 and 2.23 dB and reduces latent temporal
difference errors by \(23.44\%\) and \(21.24\%\), respectively. On the
4-step distilled model, TAA provides a smaller but consistent
\(5.03\%\) endpoint-error reduction. All three comparisons have a
10/0 paired win rate and two-sided exact sign-test \(p=0.00195\); full
confidence intervals are reported in Supplementary Table S2.

Alignment-intensity validation reveals a non-monotonic dependence on
correction strength. An intermediate intensity of 0.75 provides the
lowest endpoint error and the strongest downstream quality, whereas
weaker or full-strength corrections are less effective. This behavior
supports a calibrated rather than unconstrained trajectory correction.

\paragraph{Factorial effects.}
The \(2\times2\) analysis in Table~\ref{tab:factorial} shows that CLL
provides the principal broad quality gain, improving VBench-5 by
approximately \(0.031\)--\(0.041\) across the three evaluated
pipelines. TAA serves a more targeted role. With CLL fixed, it improves
the imaging-quality dimension by \(+0.00599\), \(+0.00254\), and
\(+0.00479\) at step 40, step 45, and on the 4-step sampler,
respectively. These detail-oriented improvements are distributed across
only one component of the five-dimensional aggregate, and therefore
appear as smaller changes in the averaged VBench-5 score.

\begin{table*}[!t]
\centering
\footnotesize
\setlength{\tabcolsep}{4.5pt}
\begin{tabular}{lrrrrrrr}
\toprule
Sampler / handoff & Trilinear & TAA+Tri. & CLL-only & TALH &
CLL gain & TAA gain & Total gain \\
\midrule
Wan50 / Step 40 (TALH-Q) & 0.77812 & 0.77901 & 0.80896 &
\textbf{0.80983} & +0.03085 & +0.00087 & +0.03171 \\
Wan50 / Step 45 (TALH-E) & 0.76776 & 0.76734 &
\textbf{0.80842} & 0.80792 & +0.04066 & -0.00050 & +0.04016 \\
Distill4 / Step 3-of-4 & 0.81796 & 0.81909 & 0.85670 &
\textbf{0.85680} & +0.03874 & +0.00010 & +0.03884 \\
\bottomrule
\end{tabular}
\caption{End-to-end factorial results on VBench-5 (\(n=10\)).
``TAA gain'' denotes the marginal effect of enabling TAA with CLL
fixed.}
\label{tab:factorial}
\end{table*}

The results support a clear division of labor. CLL restores the
cross-resolution latent representation and supplies the dominant
improvements in motion smoothness, aesthetics, and overall imaging
quality. TAA further calibrates the handoff state toward the endpoint
trajectory, enriching localized textures and fine structures that may
be underrepresented by aggregate automatic metrics.

\subsection{Perceptual Effects and Distilled-Sampler Compatibility}

Human evaluation confirms that the detail recovery introduced by TAA is
perceptually visible. With the same CLL operator, TAA receives
fine-detail win/loss/tie counts of \(8/2/0\) at step 40 and \(9/0/1\)
at step 45. Overall-quality judgments are also positive, reaching
\(6/2/2\) and \(6/0/4\), respectively. These preferences support the
dimension-level VBench observation that TAA selectively improves
imaging quality and localized structure.

The additional fine-scale content introduces a spatial--temporal
trade-off: temporal stability favors the unaligned CLL state in both
settings. This suggests that TAA shifts the operating point toward
sharper local rendering, while CLL and the remaining high-resolution
suffix provide the broader spatial and temporal reconstruction.

\begin{table}[!t]
\centering
\footnotesize
\setlength{\tabcolsep}{2pt}
\begin{tabular}{lrrr}
\toprule
Comparison & Details & Overall & Temporal \\
\midrule
TAA vs.\ Unaligned, CLL (S40) & 8/2/0 & 6/2/2 & 1/9/0 \\
TAA vs.\ Unaligned, CLL (S45) &
\textbf{9/0/1} & \textbf{6/0/4} & 0/8/2 \\
CLL-only vs.\ Trilinear (S45) &
\textbf{10/0/0} & \textbf{10/0/0} & \textbf{8/0/2} \\
TALH-E vs.\ Trilinear (S45) &
\textbf{10/0/0} & \textbf{10/0/0} & \textbf{9/0/1} \\
\bottomrule
\end{tabular}
\caption{Prompt-level majority votes (W/L/T) over ten prompts with three
raters per prompt. Bold entries have exact sign-test \(p<0.05\).}
\label{tab:human}
\end{table}

CLL provides the broader perceptual improvement: at step 45, both
CLL-only and TALH-E obtain \(10/10\) prompt-level wins over trilinear
handoff for fine details and overall quality, while also improving
temporal stability. TAA adds a narrower refinement of local detail on
top of this cross-resolution reconstruction.

TALH is also compatible with the 4-step distilled Wan sampler. On the
\(368\text{p}\) trajectory, Trilinear Handoff achieves a VBench-5 score
of 0.81796, whereas TALH-D4 reaches 0.85680. The associated cold-start
latency increases only from 167.09 to 170.29 seconds. As in the 50-step
setting, CLL supplies the primary cross-resolution quality improvement,
while TAA performs localized endpoint calibration and improves the
imaging-quality component.

Qualitative comparisons in Figure~\ref{fig:teaser}(a) further illustrate
CLL's recovery of high-frequency textures and TAA's localized
adjustments. Extended spatial crops and temporal sequences are provided
in Figure S1. The content-matched
\texttt{Native-HR (estimated)} reference is produced from a complete
\(368\text{p}\) trajectory and is used only for frame-level structural
comparison; true \(720\text{p}\) Native-HR Sampling is used for the
quantitative evaluation in Table~\ref{tab:system}.

\section{Discussion and Limitations}

TALH assigns trajectory correction, cross-resolution lifting, and
high-resolution refinement to TAA, CLL, and HTR. This decomposition
keeps CLL in a stable clean-latent domain while allowing the frozen
generator to recover high-frequency details after re-entry. Endpoint
metrics, imaging-quality gains, and human preferences indicate that TAA
improves local detail and overall perceptual quality, although its effect
depends on the CLL input distribution and high-resolution suffix.

TAA and CLL use model-endogenous self-distillation: the frozen Wan model
supplies high-resolution references for CLL and complete low-resolution
trajectories for TAA, avoiding external paired data, standalone
upscalers, and additional teachers.

Several limitations remain. CLL is tied to the target VAE, generator
distribution, and scaling configuration, so cross-model transfer is
unverified. A distribution gap persists between its training inputs and
native low-resolution trajectory states. Fixed handoff steps also
require step-specific TAA parameters, motivating timestep-conditioned
adapters for adaptive handoffs. Finally, TALH addresses in-trajectory
spatial scaling rather than real-world degradations such as compression,
blur, or sensor noise. It inherits the base generator's content risks
and safety requirements.

\section{Conclusion}

We presented TALH, a lightweight framework combining a low-resolution
prefix with a short high-resolution suffix. TAA aligns the residual
trajectory, CLL lifts clean latents across resolutions, and HTR
re-enters the frozen high-resolution trajectory, with supervision
generated entirely by the base model.

TALH-Q and TALH-E achieve \(1.83\times\) and \(2.22\times\) speedups
while retaining \(97.76\%\) and \(97.53\%\) of Native-HR VBench-5
performance. CLL consistently outperforms trilinear interpolation,
while TAA improves residual alignment and localized visual detail across
standard and distilled samplers.
